% Generated draft: requires mathematical and source review.
\documentclass[11pt]{article}
\usepackage[margin=25mm]{geometry}
\usepackage{fontspec}
\setmainfont{Latin Modern Roman}
\usepackage{amsmath,amssymb,mathtools,unicode-math}
\setmathfont{Latin Modern Math}
\usepackage{enumitem,needspace,xurl,hyperref}
\setlength{\parindent}{0pt}
\setlength{\parskip}{5pt}
\setlength{\emergencystretch}{3em}
\newcommand{\N}{\mathbb{N}}
\newcommand{\Z}{\mathbb{Z}}
\newcommand{\Q}{\mathbb{Q}}
\newcommand{\R}{\mathbb{R}}
\newcommand{\C}{\mathbb{C}}
\newcommand{\F}{\mathbb{F}}
\newcommand{\A}{\mathbb{A}}
\newcommand{\PP}{\mathbb{P}}
\newcommand{\E}{\mathbb{E}}
\newcommand{\eps}{\varepsilon}
\newcommand{\dd}{\,\mathrm{d}}
\begin{document}
\section*{UnsolvedMath ID: 1200021}
\textbf{Source problem code:} catalogue entry
\subsection{Definitions and mathematical statement}
Let \(\Gamma\) be a discrete group with Kazhdan property \((T)\): every unitary representation admitting almost invariant unit vectors has a nonzero invariant vector. In the usual bounded-degree setting, fix a finite symmetric generating set \(S\) and its labeled Cayley graph. A sofic approximation \((G_n)\) consists of finite \(S\)-labeled graphs whose rooted radius-\(r\) neighborhoods agree with that Cayley graph at a proportion tending to \(1\), for each fixed \(r\). The question asks whether there exist graphs \(G'_n\), at asymptotically vanishing edit distance from \(G_n\), such that every sequence of connected components selected along a subsequence of the indices forms an expander family.

For a finite graph \(C\) with at least two vertices, define its edge expansion by
\[
h(C)=\min_{\varnothing\ne A\subseteq V(C),\ |A|\le |V(C)|/2}
\frac{|\partial_E A|}{|A|},
\]
where \(\partial_E A\) consists of edges with exactly one endpoint in \(A\). In the standard edge-edit interpretation, the sought changes would keep the vertex sets and satisfy
\[
\frac{|E(G_n)\mathbin{\triangle}E(G'_n)|}{|V(G_n)|}\longrightarrow0.
\]
A natural uniform-expansion reading of the component conclusion is that some \(c>0\) satisfies \(h(C)\ge c\) for every nontrivial connected component \(C\) of every sufficiently late \(G'_n\). Equivalently, arbitrary component sequences along increasing indices have uniformly positive expansion, with bounded degree inherited from an appropriate modification convention. These formulas are conditional conventions: the record does not specify whether edits must preserve labels and degree, whether vertex deletion is allowed, or how edit distance is normalized.

There is also a distinction between positive expansion and the usual requirement that an expander \emph{family} have orders tending to infinity. The source does not state that \(\Gamma\) is infinite or explain bounded components and isolated vertices. If growth of component orders is required, the quantifier must specify which component subsequences qualify. No such extra restriction is silently imposed. The supplied research summary records partial progress on property-\((T)\) expansion of sofic approximations, but says this exact edit-distance statement was not verified as fully resolved.
\subsection{Short English statement}
Can a sofic approximation of a property-\((T)\) group be modified by a negligible number of edits so that all its components have expander behavior? The edit model, uniformity, and treatment of bounded components must be fixed.
\subsection{Sources}
\begin{itemize}
\item Local UnsolvedMath registry, permanent ID \(1200021\), exact source code \texttt{catalogue entry}, ``Some Questions --- Question 21.''
\item Abert, \emph{Some questions} (2010), Question 21; \url{https://www.renyi.hu/~abert/questions.pdf}; manual transcription; recorded access 2026-07-29.
\item Supplied partial-progress summary names Ab\'ert--Elek--Nikolov--Szegedy as examples of related work but supplies no complete authoritative citation or verified resolution of the exact statement.
\end{itemize}
\end{document}
